%Chargement des paquets
\usepackage{amsmath}
\usepackage{amsfonts}
\usepackage{amssymb}
\usepackage{enumerate}
\usepackage{mathtools}
\usepackage{bbm}
\usepackage{xparse, etoolbox}
\usepackage{enumerate}
\usepackage{enumitem}
\usepackage{mathabx}
\usepackage{babel}[french]

%environment
\usepackage{amsthm}
\usepackage{keytheorems}
\usepackage{hyperref} 

%Interface théorème
\newkeytheorem{lem}[
	name = Lemme
]

\newkeytheorem{prop}[
	name = Proposition
]

\newkeytheorem{thm}[
	name = Théorème
]

\newkeytheorem{coro}[
	name = Corollaire
]

\renewcommand*{\proofname}{Démonstration}

\theoremstyle{definition}
\newtheorem{defn}{Définition}

\theoremstyle{definition}
\newtheorem*{nota}{Notation}

\theoremstyle{definition}
\newtheorem*{limi}{Liminaire}

\theoremstyle{remark}
\newtheorem{rmq}{Remarque}

\theoremstyle{plain}
\newtheorem*{fait}{Fait}


%Ensembles de nombres
\newcommand{\N} {\mathbb{N}}
\newcommand{\Ne}{\N^\ast}
\newcommand{\Z} {\mathbb{Z}}
\newcommand{\Q} {\mathbb{Q}}
\newcommand{\R} {\mathbb{R}}
\newcommand{\Rb}{\overline{\mathbb{R}}}
\newcommand{\Rp}{\R_+}
\newcommand{\Rm}{\R_-}
\newcommand{\K} {\mathbb{K}}
\newcommand{\Ket} {\mathbb{K^{\ast}}}
\newcommand{\Cx}{\mathbb{C}}

%Ensembles, tribus
\newcommand{\A}{\mathbb{A}}
\renewcommand{\P}{\mathcal{P}}
\newcommand{\T}{\mathcal{T}}
\newcommand{\F}{\mathcal{F}}
\newcommand{\C} {\mathcal{C}}
\newcommand{\ouv}{\mathcal{O}}
\newcommand{\B}{\mathcal{B}}
\newcommand{\M}{\mathcal{M}}
\newcommand{\etag}{\mathcal{E}}
\newcommand{\etagOp}{\etag^{+}(\Omega)}
\newcommand{\etagO}{\etag(\Omega)}
%%Lp avant quotientage
\newcommand{\Lic}{\mathcal{L}^1}
\newcommand{\Lpc}{\mathcal{L}^p}
\newcommand{\Linfc}{\mathcal{L}^{\infty}}
\newcommand{\Lifull}{\Lic(\Omega, \B, m)}
\newcommand{\Lpfull}{\Lpc(\Omega, \B, m)}
\newcommand{\Linffull}{\Linfc(\Omega, \B, m)}
%%Lp après quotientage
\newcommand{\Li}{L^1}
\newcommand{\Ld}{L^2}
\newcommand{\Lp}{L^p}
\newcommand{\Linf}{L^{\infty}}
\newcommand{\Listd}{\Li(\Omega, m)} 
\newcommand{\Ldstd}{\Ld(\Omega, m)} 
\newcommand{\Lpstd}{\Lp(\Omega, m)} 

%Quelques opérateurs
\DeclareMathOperator{\codim}{codim}
\DeclareMathOperator{\rg}{rg}
\DeclareMathOperator{\tr}{tr}
\DeclareMathOperator{\vect}{Vect}
\DeclareMathOperator{\ima}{Im}
\DeclareMathOperator{\id}{Id}
\DeclareMathOperator{\sign}{sign}
\DeclareMathOperator{\ext}{ext}
\newcommand{\ind}{\mathbbm{1}}
\newcommand*{\spr}[2]{\left(\,#1\,\middle|\,#2\,\right)}
\newcommand{\interior}[1]{%
  {\kern0pt#1}^{\mathrm{o}}%
}
\DeclareMathOperator{\card}{card}
\DeclareMathOperator{\ppcm}{ppcm}

%Commandes de pure flemme
\newcommand{\veps}{\varepsilon}
\newcommand{\intab}{\int_{a}^{b}}
\newcommand{\intR}{\int_{-\infty}^{\infty}}
\newcommand{\intinf}{\int_{- \infty}^{+ \infty}}
\newcommand{\equi}{\Leftrightarrow}
\newcommand{\Kev}{$\K$-espace vectoriel }
\newcommand{\Kevs}{$\K$-espaces vectoriels }
\newcommand{\Rev}{$\R$-espace vectoriel }
\newcommand{\Revs}{$\R$-espaces vectoriels }
\newcommand{\Cev}{$\Cx$-espace vectoriel }
\newcommand{\Cevs}{$\Cx$-espaces vectoriels }
\newcommand{\evn}{(E, \norm{.})}
\newcommand{\interf}[2]{[#1,#2]}
\newcommand{\limb}{\lim\limits_}
\newcommand{\lebn}{\lambda_n}
\newcommand{\pp}{\text{~presque partout}}
\newcommand{\derivp}[2]{\frac{\partial #1}{\partial #2}}
\newcommand{\famiu}{(u_i)_{i \in I}}
\newcommand{\sev}{\text{~s.e.v.}}
\newcommand{\Mnp}{M_{n,p}(\K)}
\newcommand{\Mn}{M_{n}(\K)}
\newcommand{\tq}{\text{~t.q.~}}
\newcommand{\mq}{~montrer~que~}
\newcommand{\et}{\quad \text{et} \quad}
\newcommand{\ou}{\text{~ou~}}
\newcommand{\Kx}{\K[X]}


%Commandes de gars chelou
\renewcommand{\subset}{\subseteq}

%Commande restriction, ty duckduckgo
\def\restr#1#2{\mathchoice
              {\setbox1\hbox{${\displaystyle #1}_{\scriptstyle #2}$}
              \restraux{#1}{#2}}
              {\setbox1\hbox{${\textstyle #1}_{\scriptstyle #2}$}
              \restraux{#1}{#2}}
              {\setbox1\hbox{${\scriptstyle #1}_{\scriptscriptstyle #2}$}
              \restraux{#1}{#2}}
              {\setbox1\hbox{${\scriptscriptstyle #1}_{\scriptscriptstyle #2}$}
              \restraux{#1}{#2}}}
\def\restraux#1#2{{#1\,\smash{\vrule height .8\ht1 depth .85\dp1}}_{\,#2}} 

%Quelques normes
\DeclarePairedDelimiter\abs{\lvert}{\rvert}
\DeclarePairedDelimiter\labs{\bigl\lvert}{\bigr\rvert}
\DeclarePairedDelimiter\norm{\lVert}{\rVert}
\DeclarePairedDelimiterXPP\normi[1]{}\lVert\rVert{_1}{\ifblank{#1}{\:\cdot\:}{#1}}
\DeclarePairedDelimiterXPP\normd[1]{}\lVert\rVert{_2}{\ifblank{#1}{\:\cdot\:}{#1}}
\DeclarePairedDelimiterXPP\normp[1]{}\lVert\rVert{_p}{\ifblank{#1}{\:\cdot\:}{#1}}
\DeclarePairedDelimiterXPP\normq[1]{}\lVert\rVert{_q}{\ifblank{#1}{\:\cdot\:}{#1}}
\DeclarePairedDelimiterXPP\norminf[1]{}\lVert\rVert{_\infty}{\ifblank{#1}{\:\cdot\:}{#1}}

%Bracket en angle
\DeclarePairedDelimiter\angl{\langle}{\rangle}

%Set, parenthèses
\newcommand{\set}[1]{\{ #1 \}}
\newcommand{\bigset}[1]{\bigl\{ #1 \bigr\}}
\newcommand{\Bigset}[1]{\Bigl\{ #1 \Bigr\}}
\newcommand{\bigpar}[1]{\bigl( #1 \bigr)}
\newcommand{\Bigpar}[1]{\Bigl( #1 \Bigr)}

%Opitmisation des opérateurs de comparaison
\DeclarePairedDelimiter\tio{o (}{)}
\DeclarePairedDelimiter\gro{O (}{)}

\newcommand{\ort}{^{\perp}}

\pagestyle{fancy}

\fancyhead[C]{Le Théorème de projection et ses conséquences}
\fancyhead[R]{}

\begin{document}

\underline{Source :} Li - Analyse fonctionnelle - page 32.

\underline{Leçons :}
\begin{itemize}
\item  205 : Espaces complets. Exemples et applications.
\item  213 : Espaces de Hilbert. Exemples d’applications.
\item  208 : Espaces vectoriels normés, applications linéaires continues. Exemples.
\item  219 : Extremums : existence, caractérisation, recherche. Exemples et applications.
\item  253 : Utilisation de la notion de convexité en analyse.

\end{itemize}

\begin{thm}
Soit $H$ un espace de Hilbert et soit $C$ une partie convexe et fermée, non vide, de $H$. 
Alors, pour tout $x \in H$, il existe un unique $y \in C$ tel que :
\[ \norm{x-y} = d(x,C) = \inf_{c \in C} \norm{x-c} . \]
On dit que $y=P_C(x)$ est la projection de $x$ sur $C$. Cette projection est caractérisée par :
\[ y \in C \et \Re{\angl{x-y, z-y} \leq 0}, \forall z \in C . \]
\end{thm}

\begin{proof}
On démontre dans un premier temps l'existence d'un tel projeté.
Notons $d=d(x,C)=\inf_{z \in C} \norm{x-z}$.
Si $d=0$ alors $x\in \overline{C} = C$, et on a démontré l'existence du projeté (et même son unicité).
Supposons donc $d >0$. 
Par définition de la borne inférieure, pour tout $n \geq 1$, il existe $z_n \in C$ tel que :
\[ d^2 \leq \norm{x-z_n}^2 \leq d^2 + \dfrac1{n} . \]
Soient $n,p \geq 1$, en appliquant l'identité du parallélogramme à $u=x-z_n$ et $v=x-z_p$ on otient :
\[ 4 \norm*{x- \dfrac{z_n+z_p}2}^2 + \norm{z_n-z_p}^2 = 2 ( \norm{x-z_n}^2 + \norm{x-z_p}^2 ) . \]
Comme $C$ est convexe, le point $\dfrac{z_n+z_p}2$ est dans $C$, donc :
\[ \norm*{x- \dfrac{z_n+z_p}2}^2 \geq d^2 , \]
en substituant, on obtient :
\[ \norm{z_n-z_p}^2 \leq 2 \left ( d^2 + \dfrac1{n} + d^2 + \dfrac1{p} \right ) - 4d^2 = 2 \left ( \dfrac1{n} + \dfrac1{p} \right ) . \]
On en déduit que la suite $(z_n)_{n \geq 1} \subset C$ est de Cauchy dans $H$ complet, elle est donc convergente. 
On note $y$ sa limite, par fermeture de $C, y \in C$.
Enfin, le fait que :
\[ d^2 \leq \norm{x-z_n}^2 \leq d^2 + \dfrac1{n} . \]
donne, par passage à la limite $\norm{x-y}$. \\

Démontrons maintenant l'unicité de ce projeté. Notons $y, z \in C$ tels que $\norm{x-y}=\norm{x-z}=d$.
En appliquant l'identité du parallélogramme, on trouve :
\[ \norm{y-z}^2 + 4 \norm*{x - \dfrac{y-z}2}^2 = 4d^2 , \]
Or $\dfrac{y-z}2 \in C$ donc :
\[ \norm{y-z}^2 + 4d^2 \leq 4d^2 ,\]
soit $\norm{y-z}^2 \leq 0$ ce qui n'est possible que si $y=z$. \\

Il reste à démontrer la caractérisation du projeté.
Prenons $z \in C$ et $y=P_C(x)$, le segment $[y,c]$ est inclus tout entier dans $C$ donc :
\[ \forall t \in [0,1], \quad tz + 1(-t)y \in C , \]
autrement dit :
\[ \forall t \in [0,1], \quad \norm{x - tz - (1-t)y}^2 = \norm{x-y +t(y-z)}^2 \geq \norm{x-y}^2 , \]
soit en développant :
\[ \forall t \in [0,1], \quad t^2 \norm{y-z}^2 +2t \Re{\angl{x-y,y-z}} \geq 0 . \]
Cette assertion est vraie si, et seulement si, la droite $t \mapsto \norm{y-z}^2 t + 2 \Re{\angl{x-y,y-z}}$ a une ordonnée à l'origine
positive ou nulle. Or :
\[ \Re{\angl{x-y,y-z}} \geq 0 \iff \Re{\angl{x-y,z-y}} \leq 0 , \]
ce qui est à la caractérisation recherchée. \\
Réciproquement, si $y$ vérifie la caractérisation, pour tout $z \in C$, on a :
\begin{align*}
\norm{y-z}^2 & = \norm{(x-y)+(y-z)}^2 = \norm{x-y}^2 + \norm{y-z}^2 + 2 \Re{\angl{x-y,y-z}} \\
& = \norm{x-y}^2 + \norm{y-z}^2 - \underbrace{2 \Re{\angl{x-y,z-y}}}_{\leq 0} \\
& \geq \norm{x-y}^2 ,
\end{align*}
et donc $y=P_C(x)$ par unité.
\end{proof}

\begin{prop}
L'application $P_C : H \to C$ est 1-lipschitzienne, en particulier continue.
\end{prop}

\begin{proof}
Prenons $x_1, x_2 \in C, x_1 \neq x_2$ et notons $y_1, y_2$ leurs projetés respectifs. On a alors :
\[ \left \{ \begin{array}{c c c}
\Re{\angl{x_1-y_1, z-y_1}} & \leq 0 & \forall z \in C \\
\Re{\angl{x_2-y_2, z'-y_2}} & \leq 0 & \forall z' \in C
\end{array} \right. \]
On prend $z=y_2$ et $z'=y_1$ et on additionne les deux inégalités, on obtient :
\[ \Re{\angl{(x_1-y_1)-(x_2-y_2),y_2-y_1}} \leq 0 . \]
En développant $\norm{y_1-y_2}^2$, on obtient :
\begin{align*}
\norm{y_1-y_2}^2 = \Re{\norm{y_1-y_2}^2} = \Re{\angl{(y_2-x_2)+(x_2-x_1)+(x_1-y_1),y_2-y_1}} \\
& = \underbrace{\Re{\angl{(x_1-y_1)-(x_2-y_2),y_2-y_1}}}_{\leq 0} + \Re{\angl{x_2-x_1,y_2-y_1}} \\
& \leq \Re{\angl{x_2-x_1,y_2-y_1}} \\
& \leq \abs{\angl{x_2-x_1,y_2-y_1}} \leq \norm{x_2-x_1} \norm{y_2-y_1} ,
\end{align*}
d'après l'inégalité de Schwarz. Si $y_1 \neq y_2$ on obtient en divisant par $\norm{y_2-y_1}$ :
\[ \norm{P_C(x_2)-P_C(x_1)} \leq \norm{x_2 - x_1} , \]
et l'inégalité reste vraie si $y_1=y_2$.
\end{proof}

\begin{rmq}
Si l'espace sur lequel projeter est un \sev fermé de $H$, les propriétés sont plus fortes.
\end{rmq}

\begin{thm}
Si $F$ est un \sev fermé de $H$, alors l'application $P_F :H \to F$ est une application linéaire continue,
et $P_F(x)$ est l'unique point $y \in F$ tel que :
\[ y \in F \et x-y \in F\ort . \]
\end{thm}

\begin{proof}
Démontrons tout d'abord la caractérisation de $y$.
Si $y = P_F(x)$, en utilisant la caractérisation valable pour $F$ convexe et fermé :
\[ y \in F \et \Re{\angl{x-y, z-y} \leq 0}, \forall z \in F , \]
comme de plus $F$ est un \sev, on a :
\[ \forall \lambda \in \K, \forall z \in F, \quad \omega = y + \lambda z \in F . \]
Si $H$ est un espace de Hilbert réel, on a :
\[ \forall \lambda \in \R, \forall z \in F, \quad \lambda \angl{x-y, z} = \angl{x-y, \omega - y} \leq 0 , \]
ce qui entraîne nécessairement $\forall z \in F, \angl{x-y, z} = 0$ soit $x-y \in F\ort$.
Si $H$ est complexe, on a :
\[ \forall \lambda \in \R, \forall z \in F, \quad \lambda \Re{\angl{x-y, z}} = \Re{\angl{x-y, \omega - y}} \leq 0 , \]
donc $\Re{\angl{x-y, z}} = 0$. Puis, en prenant $\omega = y + i \lambda z$ :
\[ \forall \lambda \in \R, \forall z \in F, \quad \lambda \Im{\angl{x-y, z}} = \Re{\angl{x-y, \omega - y}} \leq 0 , \]
donc $\Im{\angl{x-y, z}} = 0$. 
Ainsi, dans les deux cas, on a bien $x-y \in F\ort$. \\

Réciproquement, si $y$ vérifie la caratérisation de ce théorème, on a alors :
\[ d(x,F)^2 = \inf_{z \in F} \norm{x-z}^2 = \inf_{z \in F} \norm{x-y+y-z}^2 = \inf_{z \in F} ( \norm{x-y}^2 + \norm{y-z}^2 ) 
	= \norm{x-y}^2 \]
et donc $y = P_F(x)$ par unicité. \\

Il reste à démontrer la linéarité de l'application $P_F$. 
Soient $x_1, x_2 \in H$ et $y_1, y_2$ leurs projetés respectifs. 
Soit $\lambda \in \K$, on a $\lambda y_1 + y_2 \in F$ et $\lambda (x_1 - y_1) + (x_2-y_2) \in F\ort$ ($F$ et $F\ort$ sont des \sev) 
donc $\lambda y_1 + y_2 = P_F(\lambda x_1 + x_2)$ d'après la caractérisation.
\end{proof}

\end{document}